% ============================================================================
% CHAPTER 3 - METHODS / PARAMETERS
% File: docs/thesis/methods/group6_land_eligibility.tex
% Group 6: riparian buffers, settlements, terrain, elevation, minimum area (Stage F3).
% Theory: docs/thesis/theory/ch2_land_eligibility_constraints.tex
% Values must match config/parameters.json. Lines marked % verify are open.
% ============================================================================

\section{Land Eligibility Parameters (Stage F3)}
\label{sec:params_group6}

Stage F3 evaluates eligibility on the 0.01$^\circ$ grid as the eligible
fraction of each pixel and aggregates it to the 0.05$^\circ$ cells. It is
run under three scenarios, restrictive, central and permissive. This
section defines the parameters that differ between them: the riparian
buffer, the settlement setback, slope, the elevation limit for wind and
the minimum eligible area. A legal rule is applied identically in the
three scenarios; where the law leaves no choice, the scenarios vary the
uncertain estimate to which the rule is applied. The residual land-use
factor of Section~\ref{sec:params_group1} applies only to what these
exclusions do not represent and does not repeat them.

\subsection{Riparian Buffers}
\label{subsec:g6_riparian}

Rivers come from HydroRIVERS, which includes reaches with a catchment of
at least 10~km$^2$ or a mean discharge of at least 0.1~m$^3$/s
\cite{lehner2013}. The channel width of each reach is estimated from its
long-term mean discharge $Q$ with the global hydraulic-geometry relation
of Andreadis et al.\ \cite{andreadis2013},
\begin{equation}
    W = 7.2\,Q^{\,b}, \qquad b = 0.48,\ 0.50,\ 0.52
    \label{eq:g6_width}
\end{equation}
in the permissive, central and restrictive scenarios, the published
uncertainty of the exponent. Because the relation was fitted to bankfull
discharge, which is larger than mean discharge, $W$ is a width proxy that
tends to underestimate the legal channel width; the bias is toward the
permissive side and is recorded as a limitation.

Each reach excludes a corridor of half-width $W/2 + B_c(W)$ around its
centre line, where $B_c$ is the buffer of country $c$ measured from the
channel edge (Table~\ref{tab:g6_riparian}). The corridor is combined with
the water classes of the land-cover map by union, so the channel is not
counted twice.

In Brazil, $B$ follows the five width classes of the Forest Code
\cite{brasil2012forestcode}. The law excludes ephemeral channels, which
HydroRIVERS does not distinguish; all reaches are treated as perennial or
intermittent. In Portugal, the margin of the public water domain is
10~m for non-navigable waters, 30~m for navigable or floatable waters and
public reservoirs, and 50~m for waters under maritime and port
jurisdiction \cite{portugal2005lei54}. Navigability is not available in
global hydrography. The central scenario uses the estimated width as a
proxy, with 10~m for reaches narrower than 10~m and 30~m for wider
reaches (Tier 3); the permissive scenario applies 10~m and the
restrictive scenario 30~m to all reaches. The 50~m margin of estuaries
is not represented, since the tidal limit is not known in the river
network. In India no national statutory riparian buffer was identified;
the buffers are author judgment (Tier 3), chosen within the range used
in siting studies.

\begin{table}[htbp]
    \centering
    \caption{Riparian buffer $B$ measured from the channel edge.}
    \label{tab:g6_riparian}
    \small
    \begin{tabular}{llccc}
        \toprule
        Country & Channel width & Restrictive & Central & Permissive \\
        \midrule
        BRA & $W<10$~m          & 30~m  & 30~m  & 30~m \\
            & $10\le W<50$~m    & 50~m  & 50~m  & 50~m \\
            & $50\le W<200$~m   & 100~m & 100~m & 100~m \\
            & $200\le W<600$~m  & 200~m & 200~m & 200~m \\
            & $W\ge 600$~m      & 500~m & 500~m & 500~m \\
        PRT & $W<10$~m          & 30~m  & 10~m  & 10~m \\
            & $W\ge 10$~m       & 30~m  & 30~m  & 10~m \\
        IND & all reaches       & 200~m & 100~m & 30~m \\
        \bottomrule
    \end{tabular}
\end{table}

In Brazil the scenarios differ through the exponent $b$, which moves
reaches between width classes. With $b=0.50$, the class limits of 10, 50,
200 and 600~m correspond to mean discharges of about 1.9, 48, 770 and
6\,900~m$^3$/s.

In Portugal, the National Ecological Reserve (REN), which protects
flood-prone areas, river beds and margins, steep slopes and coastal
systems, is not represented in any scenario. Its regime admits renewable
plants in several of its typologies under prior communication or
authorization \cite{portugal2008ren}, so it is not a general prohibition.
% verify: admissible uses of renewable plants per REN typology (RJREN annex)
The vector layer published by the Directorate-General for the Territory
was not available in complete form when the analysis was run, and a
partial layer was not used, because it would exclude some municipalities
and not others.

\subsection{Settlements}
\label{subsec:g6_settlements}

Settlements are taken from the GHS settlement model, which applies the
Degree of Urbanisation to gridded population and built-up surface
\cite{schiavina2023smod}. A settlement is any 1~km cell classified as
rural cluster (village) or denser, that is, the classes from rural
cluster to urban centre. Dispersed rural population, the low-density and
very-low-density rural classes, is not excluded by this criterion. The
settlement mask is the same in the three scenarios; the scenarios vary
the setback around it (Table~\ref{tab:g6_setback}).

\begin{table}[htbp]
    \centering
    \caption{Setback from settlements.}
    \label{tab:g6_setback}
    \small
    \begin{tabular}{lcccl}
        \toprule
        Technology & Restrictive & Central & Permissive & Source \\
        \midrule
        PV   & 1\,000~m & 500~m & 0~m   & \cite{uyan2013,watson2015}; restrictive Tier 3 \\
        Wind & 1\,000~m & 800~m & 500~m & \cite{rediske2021,drechsler2017,lehmann2021,watson2015} \\
        \bottomrule
    \end{tabular}
\end{table}

Settlements are not defined by a population-density threshold. A density
threshold applied uniformly would remove very different shares of land in
countries with different rural settlement patterns, while the settlement
classes combine density with contiguity and minimum population. A
threshold of 200 inhabitants/km$^2$ is run as a sensitivity.

\subsection{Terrain and Elevation}
\label{subsec:g6_terrain}

Slope is computed in percent at 30~m from FABDEM, a version of the
Copernicus GLO-30 model with forests and buildings removed
\cite{hawker2022fabdem}, and a pixel's eligible fraction excludes the
share of its area steeper than the threshold. GLO-30 is a surface model:
vegetation and buildings add relief that, at a 30~m spacing, exceeds the
low slope limits of PV. Over the three countries, the PV-eligible area
obtained with FABDEM is 1.12 to 1.34 times the area obtained with GLO-30
at the central limit and 1.24 to 1.63 times at the restrictive limit.
For wind the two models differ by 0.1 to 2.1\,\% at the central and
permissive limits; at the restrictive limit of 10\,\% the FABDEM area is
larger by 4.2\,\% in Brazil, 1.6\,\% in India and 11.1\,\% in Portugal.
The two results bound the sensitivity of eligibility to the elevation
model, and FABDEM is the model adopted. Elevation, used for air density and the wind
elevation limit, remains from GLO-30, where the height of vegetation is
immaterial. Elevation limits apply only to wind,
because air density already enters the wind capacity factor
(Section~\ref{sec:params_group1}); the limit stands for access,
construction and cold-climate constraints, and the permissive scenario
applies none.

\begin{table}[htbp]
    \centering
    \caption{Slope and elevation limits.}
    \label{tab:g6_terrain}
    \small
    \begin{tabular}{llcccl}
        \toprule
        Parameter & Technology & Restrictive & Central & Permissive & Source \\
        \midrule
        Maximum slope & PV   & 3\,\%  & 5\,\%  & 10\,\% & \cite{lopez2012,tao2024,blm2024westernsolar} \\
        Maximum slope & Wind & 10\,\% & 20\,\% & 30\,\% & \cite{vanhaaren2011,lopez2012,epg2024renewland} \\
        Maximum elevation & Wind & 2\,000~m & 3\,000~m & none & \cite{epg2024renewland}; central Tier 3 \\
        \bottomrule
    \end{tabular}
\end{table}

\subsection{Minimum Eligible Area}
\label{subsec:g6_minarea}

Eligible pixels are grouped into clusters of 8-connected 0.01$^\circ$
pixels with a positive eligible fraction, and a cluster is kept only if
its eligible area reaches a minimum $A_{\min}$. For PV, the central value
is 1~km$^2$ \cite{lopez2012}. For wind, the central scenario additionally
discards pixels whose eligible fraction is below 10\,\%, the share needed
for one turbine per square kilometre \cite{elsner2023onshore}. Because
the pixel area varies with latitude, 10\,\% corresponds to about 9 to
12~ha across the study area, against the 10~ha of the original rule. The
restrictive scenario sets $A_{\min}$ to the area of the 100~MW reference
project of Section~\ref{subsec:g3_boundary}, $A_{\min} =
P_{\mathrm{ref}}/\mathrm{PD}$, with the power density PD of each
technology, so that every retained cluster can host the project to which
the connection cost refers. The permissive scenario applies no minimum.

\begin{table}[htbp]
    \centering
    \caption{Minimum eligible area.}
    \label{tab:g6_minarea}
    \small
    \begin{tabular}{lccc}
        \toprule
        Technology & Restrictive & Central & Permissive \\
        \midrule
        PV   & cluster $\ge P_{\mathrm{ref}}/\mathrm{PD}_{\mathrm{PV}}$
             & cluster $\ge 1$~km$^2$ & none \\
        Wind & cluster $\ge P_{\mathrm{ref}}/\mathrm{PD}_{\mathrm{wind}}$; pixel $\ge 10\,\%$
             & pixel $\ge 10\,\%$ & none \\
        \bottomrule
    \end{tabular}
\end{table}

\subsection{Limitations}
\label{subsec:g6_limitations}

Channel widths are estimated from mean discharge with a bankfull
relation, which underestimates the legal width, and HydroRIVERS omits
small streams and does not distinguish ephemeral channels or navigable
waters. In Portugal, the REN and flood zones are not represented, so the
restrictive scenario is an upper bound with respect to them. The Indian
riparian buffers and several scenario bounds are author judgment. Settlement setbacks come from studies in other
countries, since none of the three countries sets a national setback for
these technologies in the sources reviewed. The settlement model has a
1~km resolution, so small villages are represented coarsely. Slope
thresholds depend on the resolution of the elevation model, and FABDEM
removes vegetation and buildings by a model rather than by observation. Clusters are
formed from pixels with any eligible fraction, so a cluster of
fragmented pixels can pass the area test without being contiguous at the
sub-pixel scale.